%% Lab Report: AI Project – Connect Four
%%
% !TeX encoding = utf8
% !TeX program = pdflatex
% !TeX spellcheck = en_US
% !BIB program = biber

\RequirePackage[utf8]{inputenc}
\RequirePackage{hgbpdfa}

\documentclass[language=english,titlepage=false,smartquotes=true]{hgbreport}

\renewcommand{\chapter}[1]{}
\graphicspath{{images/}}
\bibliography{hgbreferences}
\ExecuteBibliographyOptions{backref=false}

\setcounter{chapter}{1}
%%%-----------------------------------------------------------------------------
\begin{document}
%%%-----------------------------------------------------------------------------

\author{Martin Scheuchenpflug, Niklas Etzlinger}
\title{Artificial Intelligence -- SS 2026\\
       Lab Report \arabic{chapter}: AI Project Connect Four}
\date{\today}

\maketitle

\begin{abstract}\noindent
This report documents the experimental analysis of the Minimax algorithm with
alpha-beta pruning, applied to the game Connect Four. Five progressively
refined versions of the \texttt{MinimaxPlayer} are evaluated by measuring the
number of board positions examined per move at search depth~4. The
modifications include removing pruning entirely, tightening the pruning
condition to include equal values, enabling deep (full-window) alpha-beta
pruning, and applying a move-ordering heuristic that prefers centre columns.
\end{abstract}

%%%-----------------------------------------------------------------------------
\section{Measuring the Original Version}
%%%-----------------------------------------------------------------------------

The original \texttt{MinimaxPlayer} uses \emph{shallow} alpha-beta pruning:
each recursive call passes only the immediate parent's current best value as a
bound, discarding information from higher ancestors. The search depth is set
to~4 (the application default) throughout all experiments.

The number of positions examined by the computer per move was recorded for
10 consecutive moves of a representative game (Table~\ref{tab:results}).
In the original version the counts range from roughly 650 to 2400, depending
on how many cutoffs the current board state enables.

%%%-----------------------------------------------------------------------------
\section{Pure Minimax -- No Pruning}
\label{sec:q2}
%%%-----------------------------------------------------------------------------

\subsection{Code change}

Both alpha-beta conditions were commented out so that no branch is ever cut:

\begin{JavaCode}[numbers=none]
// in expandMaxNode
// FIXME: a-b prune
// if (strength > parentMinimum) {
//   board.undoLastMove();
//   return strength;
// }

// in expandMinNode
// FIXME: a-b prune
// if (strength < parentMaximum) {
//   board.undoLastMove();
//   return strength;
// }
\end{JavaCode}

\subsection{Formula for the exact move count}

With pruning disabled the algorithm visits every node in the search tree. The
board has $n = 8$ columns. Because each column holds 6 rows, all 8 columns
remain available for the first many plies, so the branching factor is a
constant~8 throughout the tree.

The root (\texttt{calculateMove}) tries 8 moves and calls
\texttt{expandMinNode}; each of those tries 8 moves and calls
\texttt{expandMaxNode}; and so on for $d$~levels in total. The total node
count is therefore the sum of a geometric series with ratio~8:

\begin{equation}
  \text{total}(d)
    = 8 + 8^{2} + 8^{3} + \cdots + 8^{d}
    = \frac{8^{d+1} - 8}{7}.
  \label{eq:formula}
\end{equation}

For the default depth $d = 4$:
\[
  \frac{8^{5} - 8}{7} = \frac{32768 - 8}{7} = \frac{32760}{7} = 4680.
\]

The measured values confirm this: the first three moves of Question~2 all
yield exactly 4680 (Table~\ref{tab:results}). Later moves fall slightly below
4680 because some search paths encounter a terminal position (four-in-a-row)
before reaching full depth and return without expanding further.

%%%-----------------------------------------------------------------------------
\section{Pruning with Equal Values}
%%%-----------------------------------------------------------------------------

\subsection{Code change}

The original pruning uses strict inequalities: a value \emph{equal} to the
parent's bound does not trigger a cutoff. Replacing \texttt{>} with
\texttt{>=} (and \texttt{<} with \texttt{<=}) prunes these tie cases as well:

\begin{JavaCode}[numbers=none]
// expandMaxNode: beta cutoff now also fires on equality
if (strength >= parentMinimum) {
    board.undoLastMove();
    return strength;
}

// expandMinNode: alpha cutoff now also fires on equality
if (strength <= parentMaximum) {
    board.undoLastMove();
    return strength;
}
\end{JavaCode}

\subsection{Effect}

The move counts decrease consistently compared to Question~1 (Table~\ref{tab:results}).
Ties are frequent early in the game: the heuristic assigns equal scores to
many symmetric positions, so the additional cutoffs on equal values can
eliminate entire subtrees that the strict version would have explored.

%%%-----------------------------------------------------------------------------
\section{Deep Alpha-Beta Pruning}
%%%-----------------------------------------------------------------------------

\subsection{Motivation}

In the original code each function passes only one bound to its child -- either
the parent MAX node's current best (alpha) or the parent MIN node's current
minimum (beta). Information from grandparents and higher ancestors is silently
discarded. \emph{Deep} alpha-beta threads both bounds through the entire tree:
every node sees the tightest window established by \emph{all} its ancestors.

\subsection{Code change}

Both recursive functions were given a full $(\alpha, \beta)$ signature and
update their respective bound locally before passing both values down:

\begin{JavaCode}[numbers=none]
private int expandMaxNode(int depth, int alpha, int beta) {
    ...
    int strength = expandMinNode(depth - 1, alpha, beta);
    if (strength >= beta) {          // beta cutoff
        board.undoLastMove();
        return strength;
    }
    if (strength > maxStrength) maxStrength = strength;
    if (strength > alpha) alpha = strength;  // tighten window
    ...
}

private int expandMinNode(int depth, int alpha, int beta) {
    ...
    int strength = expandMaxNode(depth - 1, alpha, beta);
    if (strength <= alpha) {         // alpha cutoff
        board.undoLastMove();
        return strength;
    }
    if (strength < minStrength) minStrength = strength;
    if (strength < beta) beta = strength;    // tighten window
    ...
}
\end{JavaCode}

The root call initialises the window to the widest possible range:
\begin{JavaCode}[numbers=none]
int strength = expandMinNode(depth - 1, maxStrength, MAX_POSSIBLE_STRENGTH);
\end{JavaCode}

\subsection{Effect}

Move counts drop further compared to Question~3 (Table~\ref{tab:results}).
Whenever the $(\alpha, \beta)$ window is tightened at any level, all
descendants benefit immediately rather than only the direct children.
The gain is most visible in positions where good moves appear early and
establish tight bounds quickly.

%%%-----------------------------------------------------------------------------
\section{Move-Ordering Heuristic}
%%%-----------------------------------------------------------------------------

Alpha-beta pruning is most effective when the best moves are considered first,
because early good moves narrow the $(\alpha, \beta)$ window quickly and cause
more cutoffs later. For Connect Four, centre columns are generally stronger
than edge columns. A fixed middle-out order was implemented:

\begin{JavaCode}[numbers=none]
// 8-column board: centre pair first, alternating outward
private static final int[] COLUMN_ORDER = {3, 4, 2, 5, 1, 6, 0, 7};

private Move[] reorderMoves(Move[] moves) {
    Move[] ordered = new Move[moves.length];
    for (int i = 0; i < COLUMN_ORDER.length; i++) {
        ordered[i] = moves[COLUMN_ORDER[i]];
    }
    return ordered;
}
\end{JavaCode}

\texttt{reorderMoves} is called at all three loop sites:
\texttt{calculateMove}, \texttt{expandMaxNode}, and \texttt{expandMinNode}.

\subsection{5a -- Pure Minimax with Middle-Out Ordering}

With alpha-beta pruning disabled the traversal order has no effect on the
\emph{number} of nodes visited: every node is expanded regardless. The move
counts therefore remain close to 4680 and follow the same formula as
Section~\ref{sec:q2} (Table~\ref{tab:results}). Minor deviations from 4680
appear in later moves because the computer now plays different moves (equal
heuristic values are resolved in a different order), leading to a different
board state and different numbers of full columns.

\subsection{5b -- Deep Alpha-Beta with Middle-Out Ordering}

Combining deep alpha-beta (Question~4) with middle-out ordering produces the
lowest move counts of all variants (Table~\ref{tab:results}). By trying the
best columns first, the algorithm establishes a tight $(\alpha, \beta)$ window
early and cuts off the remaining columns with minimal effort. The reduction
is dramatic -- move~1 drops from 993 (Question~4) to just 302.

%%%-----------------------------------------------------------------------------
\section*{Results Summary}
%%%-----------------------------------------------------------------------------

Table~\ref{tab:results} collects all measurements. Notes on comparability:
Questions~1--4 were measured on the same sequence of human moves (the computer
responds identically in all versions that use the same evaluation order).
Questions~5a and~5b diverge from move~2 onwards because the computer's choices
change under middle-out ordering.

\begin{table}[H]
  \centering\small
  \caption{Number of moves examined per computer turn at depth~4.
           Columns correspond to the six experimental variants.}
  \label{tab:results}
  \begin{tabular}{r rrrrrr}
    \hline
    \textbf{Move} &
    \textbf{Q1} &
    \textbf{Q2} &
    \textbf{Q3} &
    \textbf{Q4} &
    \textbf{Q5a} &
    \textbf{Q5b} \\
    & \small original &
    \small pure minimax &
    \small ${\geq}$ prune &
    \small deep a-b &
    \small pure + order &
    \small deep + order \\
    \hline
     1 & 1426 & 4680 & 1208 &  993 & 4680 & 302 \\
     2 &  925 & 4680 &  771 &  642 & 4680 & 360 \\
     3 &  654 & 4680 &  574 &  514 & 4650 & 702 \\
     4 &  809 & 4623 &  745 &  612 & 4538 & 899 \\
     5 & 1221 & 4622 & 1158 &  934 & 4650 & 872 \\
     6 & 2459 & 4056 & 2007 & 1936 & 4649 & 773 \\
     7 &  909 & 4544 &  854 &  754 & 4333 & 785 \\
     8 &  760 & 4256 &  736 &  664 & 4325 & 671 \\
     9 &  899 & 3978 &  793 &  665 & 2792 & 305 \\
    10 &  664 & 3641 &   59 &   59 & 2347 &  95 \\
    \hline
  \end{tabular}
\end{table}
\clearpage

%%%-----------------------------------------------------------------------------
\section{Reducing to 7 Columns}
%%%-----------------------------------------------------------------------------

\subsection{Changing the column count}

Standard Connect Four is played on a $6 \times 7$ board. The original code
uses \texttt{NUMBER\_OF\_COLUMNS = 8}. Changing this constant to~7 is
sufficient to play the game correctly; all array sizes and loop bounds
throughout \texttt{C4Board.java} are derived from the constant at compile time.

\begin{JavaCode}[numbers=none]
public static final int NUMBER_OF_COLUMNS = 7; // was 8
\end{JavaCode}

\subsection{The magic number 84}

One line in \texttt{C4Board.java} is not derived from the constants:

\begin{JavaCode}[numbers=none]
rows = new Vector<C4Row>(84);
\end{JavaCode}

This sets the initial capacity of the vector that stores every possible
group of four consecutive slots -- the groups used by the heuristic to
evaluate the board. The number~84 is the exact count of such groups on
a $6 \times 8$ board, computed direction by direction:

\begin{center}
\begin{tabular}{lll}
  \hline
  Direction & Formula & $6 \times 8$ \\
  \hline
  Horizontal          & $R \cdot (C-3)$         & $6 \cdot 5 = 30$ \\
  Vertical            & $(R-3) \cdot C$         & $3 \cdot 8 = 24$ \\
  Diagonal $\nearrow$ & $(R-3) \cdot (C-3)$     & $3 \cdot 5 = 15$ \\
  Diagonal $\searrow$ & $(R-3) \cdot (C-3)$     & $3 \cdot 5 = 15$ \\
  \hline
  Total               &                         & $84$ \\
  \hline
\end{tabular}
\end{center}

The general formula in terms of \texttt{NUMBER\_OF\_ROWS} ($R$) and
\texttt{NUMBER\_OF\_COLUMNS} ($C$) is:

\begin{equation}
  R(C-3) + (R-3)C + 2(R-3)(C-3).
  \label{eq:rows}
\end{equation}

For the standard $6 \times 7$ board this gives
$6 \cdot 4 + 3 \cdot 7 + 2 \cdot 3 \cdot 4 = 24 + 21 + 24 = 69$.

The fix replaces the hard-coded literal with the formula so it remains
correct for any board size:

\begin{JavaCode}[numbers=none]
rows = new Vector<C4Row>(
    NUMBER_OF_ROWS * (NUMBER_OF_COLUMNS - 3)
    + (NUMBER_OF_ROWS - 3) * NUMBER_OF_COLUMNS
    + 2 * (NUMBER_OF_ROWS - 3) * (NUMBER_OF_COLUMNS - 3));
\end{JavaCode}

The move-ordering array in \texttt{MinimaxPlayer.java} was also updated from
the 8-element sequence to the correct 7-element centre-out order:

\begin{JavaCode}[numbers=none]
// 7-column board: centre column 3 first, alternating outward
private static final int[] COLUMN_ORDER = {3, 2, 4, 1, 5, 0, 6};
\end{JavaCode}

After these three changes the game runs correctly on the standard
$6 \times 7$ board.

%%%-----------------------------------------------------------------------------
\section{The Heuristic Evaluation Function}
%%%-----------------------------------------------------------------------------

\subsection{What quantities does it compute, and how are they weighted?}

The heuristic counts \emph{open groups}: groups of 4 consecutive slots
(horizontal, vertical, or diagonal) that contain tokens of only one player---the
other player has placed no token in the group yet, so it can still be
completed into a four-in-a-row. A group with tokens of both players is
\emph{dead} and counted for neither. For each player the number of open groups
of size 1, 2, and 3 is tracked separately.

The raw strength of a player is computed in \texttt{C4Stats}:

\begin{JavaCode}[numbers=none]
// SHIFT_3_IN_A_ROW = 5  =>  x32
// SHIFT_2_IN_A_ROW = 2  =>  x4
return (p1_2InARow << SHIFT_2_IN_A_ROW)
     + (p1_3InARow << SHIFT_3_IN_A_ROW)
     + p1_1InARow;
// i.e. 32*(3-in-a-row) + 4*(2-in-a-row) + 1*(1-in-a-row)
\end{JavaCode}

The comment in the source claiming multipliers of 64 and 16 is outdated; the
actual weights are \textbf{32\,:\,4\,:\,1}. Two special cases override the
formula: a player with 4-in-a-row yields \texttt{Integer.MAX\_VALUE} (win);
if the \emph{opponent} has 4-in-a-row the raw strength is~\texttt{0}.
The value returned to minimax is always the difference:
\[
  \texttt{getStrength}(p) = \text{strength}_p - \text{strength}_{\overline{p}}.
\]

\subsection{How is it computed? Object communication}

The evaluation is never recomputed from scratch. It is maintained
\emph{incrementally} through a chain of event notifications that follow the
Observer pattern.

\textbf{C4Slot} is the \emph{subject}. It stores the content of one cell and
holds a list of \texttt{C4SlotListener} objects. Whenever \texttt{setContents()}
is called it fires \texttt{contentsChanged(old,~new)} on every registered
listener.

\textbf{C4Row} is the \emph{observer}. Each instance represents one group of
4 consecutive slots and implements \texttt{C4SlotListener}. Its constructor
registers it on all four slots. It maintains counters \texttt{numPlayer1} and
\texttt{numPlayer2}. On notification it updates these counters and calls the
matching \texttt{Inc}/\texttt{Dec} method on the shared \texttt{C4Stats}
object---for example, if player~1 now has 2 tokens in the group and the
opponent has none, it calls \texttt{stats.p1\_1InARowDec()} followed by
\texttt{stats.p1\_2InARowInc()}.

\textbf{C4Stats} is the \emph{global accumulator}. It holds board-wide
counters that \texttt{getStrength()} simply reads---evaluation costs
\emph{O(1)}. A slot can participate in up to 11 overlapping groups, which
is why \texttt{C4Slot} pre-allocates its listener array with capacity~11.

The full call chain on every \texttt{board.move()} is:

\begin{JavaCode}[numbers=none]
C4Board.move()
  -> slot.setContents(playerNumber)
    -> notifies each listening C4Row
      -> row.contentsChanged() updates numPlayer1 / numPlayer2
        -> calls stats.pX_NInARowInc() or Dec()
\end{JavaCode}

\texttt{undoLastMove()} triggers the same chain in reverse, keeping all
counters consistent without any board scan.

\subsection{The restart bug}

The counter fields in \texttt{C4Stats} are declared \texttt{static}, making
them class-level variables shared across \emph{all} instances. When
\texttt{C4Board.clone()} creates \texttt{clone.stats = new C4Stats()}, the new
object still reads and writes the same static counters as the original board.
The board-state copy loop then calls \texttt{setContents()} on the clone's
slots, which fires the clone's freshly created \texttt{C4Row} listeners and
increments the shared counters a second time. The current board state is thus
added as a constant offset to every counter.

Because the offset is identical for every leaf node in a single minimax call,
the relative ordering of moves is unaffected---the extra constant cancels in
the difference $\text{strength}_p - \text{strength}_{\overline{p}}$. The bug
is therefore harmless in practice.

%%%-----------------------------------------------------------------------------
\section{Missing the Immediate Win}
%%%-----------------------------------------------------------------------------

Sometimes the computer can win in one move but plays a different move instead.
The root cause is that the heuristic assigns \emph{exactly the same score} to
every winning position regardless of how many moves away the win is:
\texttt{firstPlayerStrength()} returns \texttt{Integer.MAX\_VALUE} whenever
\texttt{p1\_4InARow != 0}, whether the win occurred at depth~1 or depth~3.

The move-selection loop in \texttt{calculateMove} updates its best move only
on a \emph{strict} improvement:

\begin{JavaCode}[numbers=none]
if (strength > maxStrength) {
    maxStrength = strength;
    maxIndex = i;
}
\end{JavaCode}

Consider the following situation at depth~4 with left-to-right move ordering:

\begin{enumerate}
  \item Move at column~0 is tried first. The search finds a forced win in
        3~more moves somewhere in the tree and returns
        \texttt{Integer.MAX\_VALUE}. \texttt{maxIndex = 0}.
  \item Move at column~3 is the immediate one-move win. After placing the
        piece \texttt{board.isGameOver()} is true, so
        \texttt{expandMinNode} returns \texttt{Integer.MAX\_VALUE}
        immediately.
  \item \texttt{Integer.MAX\_VALUE > Integer.MAX\_VALUE} is \textbf{false},
        so \texttt{maxIndex} remains~0.
  \item The computer plays column~0---a winning move, but not the
        immediate one.
\end{enumerate}

The computer still wins the game; it simply does not do so as quickly as
possible. A fix would be to encode the win depth in the score, e.g.\
\texttt{Integer.MAX\_VALUE\,$-$\,depth}, so shallower wins always outrank
deeper ones. The current code makes no such distinction.

%%%-----------------------------------------------------------------------------
\section*{Summary and Comments}
%%%-----------------------------------------------------------------------------

The experiments clearly show the cumulative effect of each optimisation.
Removing pruning (Q2) inflates the count to the theoretical maximum of
$\frac{8^{5}-8}{7} = 4680$ nodes. Restoring shallow pruning with equal-value
cutoffs (Q3) already reduces counts substantially below the original (Q1).
Deep alpha-beta (Q4) cuts further by propagating bounds from all ancestors.
Finally, middle-out move ordering (Q5b) delivers the largest single reduction:
starting from the best columns narrows the search window immediately and
causes the remaining columns to be pruned with far less work.

The Q5a result confirms the theoretical expectation that move ordering is
irrelevant for pure minimax -- the counts stay near 4680 regardless of column
order, with small deviations caused by the diverging game state.

Question~6 shows that a single constant controls the board width throughout
the program, with the exception of the hard-coded vector capacity~84. That
literal encodes the number of four-in-a-row groups on the original
$6 \times 8$ board; its general form (Equation~\ref{eq:rows}) evaluates to~69
for the standard $6 \times 7$ configuration.

Question~7 reveals an elegant incremental design: the heuristic is maintained
at O(1) cost per move through the Observer pattern, with \texttt{C4Row}
objects listening to individual slots and propagating changes into a shared
\texttt{C4Stats} accumulator.

%%%-----------------------------------------------------------------------------
\end{document}
%%%-----------------------------------------------------------------------------
